\section{Constraint Satisfaction Problems (CSPs)}
\subsubsection*{Was sind CSPs?}
Ein \fatsf{Constraint Satisfaction Problem (CSP)} ist ein Problem, das gelöst wird, indem eine \fatsf{Zuweisung von Werten zu Variablen} gefunden wird, die bestimmte \fatsf{Einschränkungen (Constraints)} erfüllt\\
\fatsf{Bestandteile eines CSPs:}
\begin{enumerate}
    \item \fatsf{Variablen $X_{i}$:} Zustände des Problems (z.B. Farbe eines Landes auf einer Karte)
    \item \fatsf{Domänen $D_{i}$:} Mögliche Werte, die den Variablen zugewiesen werden können
    \item \fatsf{Constraints C:} Regeln, die angeben, welche Werte erlaubt sind (z.B. \enquote{Benachbarte Länder müssen unterschiedliche Farben haben})
\end{enumerate}
\fatsf{Lösung eines CSP:}\\
Eine gültige Zuweisung von Werten zu allen Variablen, die \fatsf{keinen Constraint} verletzt
\subsubsection*{Beispiele für reale CSPs}
\fatsf{Typische Anwendungen:}
\begin{itemize}
    \item \fatsf{Zuweisungsprobleme:} Wer unterrichtet welche Klasse?
    \item \fatsf{Stundenplanung:} Wann und wo wird eine Klasse angeboten?
    \item \fatsf{Hardware-Konfiguration:} Optimale Ressourcennutzung in Rechenzentren
    \item \fatsf{Terminplanung:} Automatische Planung von Meetings
    \item \fatsf{Grundrissplanung:} Optimale Raumnutzung
\end{itemize}
\fatsf{Spezialfälle:}
\begin{itemize}
    \item \fatsf{Lineare Constraints:} Lassen sich mit \fatsf{linearem Programmieren} in \fatsf{polynomiealer Zeit} lösen
    \item \fatsf{Nichtlineare Constraints:} Oft \fatsf{unentscheidbar} (keine generelle Lösungsmethode)
\end{itemize}
\subsubsection*{Arten der Variablenzuweisung}
Ein CSP-Zustand wird durch eine Zuweisung von Werten zu Variablen definiert:
\begin{enumerate}
    \item \fatsf{Konsistente (Legale) Zuweisung:}
    \begin{itemize}
        \item Eine Zuweisung, die \fatsf{keine} Constraints verletzt
    \end{itemize}
    \item \fatsf{Komplette Zuweisung:}
    \begin{itemize}
        \item Jede Variable ist zugewiesen, und die Lösung bleibt konsistent
    \end{itemize}
    \item \fatsf{Partielle Zuweisung:}
    \begin{itemize}
        \item Nur einige Variablen sind zugwiesen (wird oft in Suchalgorithmen verwendet)
    \end{itemize}
\end{enumerate}
\fatsf{Warum keine einfache Suche?}\\
Da Lösungen \fatsf{durch Constraints} definiert sind, können effizientere Techniken genutzt werden als \fatsf{Standard-Suchalgorithmen}
\begin{figure}[h]
    \centering
    \includegraphics[width=0.5\linewidth]{CSP_Images/n-Queens.jpeg}
\end{figure}
\subsubsection*{Constraint-Graphen}
Ein \fatsf{Constraint-Graph} abstrahiert das Problem und macht es leichter lösbar\\
\fatsf{Wie wird ein Constraint-Graph erstellt?}
\begin{itemize}
    \item \fatsf{Jede Variable} wird als \fatsf{Knoten} dargestellt
    \item \fatsf{Jede Einschränkung} zwischen Variablen wird als \fatsf{Kante} dargestellt
\end{itemize}
\subsubsection*{Das Vier-Farben-Theorem}
\begin{figure}[h]
    \centering
    \includegraphics[width=0.7\linewidth]{CSP_Images/The Four Color Map Theorem.jpeg}
\end{figure}
\begin{itemize}
    \item Beweis dauerte über \fatsf{100 Jahre} und wurde erstmals mit einem Computer abgeschlossen!
\end{itemize}
\subsubsection*{Beispiel: Arithmetisches CSP - TWO + TWO = FOUR}
\fatsf{Ziel:} Finde eine \fatsf{Zuweisung von Ziffern zu Buchstaben}, die die Gleichung erfüllt\\
\fatsf{Constraints:}
\begin{itemize}
    \item Jede Variable repräsentiert eine Ziffer 0 - 9
    \item Gleiche Buchstaben müssen dieselbe Ziffer haben
    \item Unterschiedliche Buchstaben müssen unterschiedliche Ziffern haben
    \item Addition muss korrekt sein
\end{itemize}
\fatsf{Technik:}
\begin{itemize}
    \item Nutzung von Constraint-Suchverfahren (Backtracking, Forward Checking, Arc Consistency)
\end{itemize}
\begin{figure}[h]
    \centering
    \includegraphics[width=1\linewidth]{CSP_Images/TWO-TWO-FOUR.jpeg}
\end{figure}
\subsubsection*{Types of Constraints in CSPs}
CSPs enthalten unterschiedliche \fatsf{Arten von Constraints}, die bestimmen, welche Werte den Variablen zugewiesen werden können
\begin{enumerate}
    \item \fatsf{Arten von Constraints}
    \begin{itemize}
        \item \fatsf{Unäre Constraints} (Einzelvariable):
        \begin{itemize}
            \item Betreffen \fatsf{nur eine Variable}
            \item Beispiel: Süd-Australien $\neq$ grün
        \end{itemize}
        \item \fatsf{Binäre Constraints} (Paarweise Einschränkungen):
        \begin{itemize}
            \item Betreffen \fatsf{zwei Variablen}
            \item Beispiel: Süd-Australien $\neq$ West-Australien (benachbarte Staaten müssen unterschiedliche Farben haben)
        \end{itemize}
        \item \fatsf{Höhere-Ordnungs-Constraints} (Drei oder mehr Variablen)
        \begin{itemize}
            \item Beispiel: Eine Gruppe von Variablen muss zusammen eine bestimmte Bedingung erfüllen
        \end{itemize}
        \item \fatsf{Weiche Constraints (Soft Constraints):}
        \begin{itemize}
            \item keine festen Regeln, aber wünschenswerte Bedingungen
            \item Beispiel: Rot ist besser als Grün (Präferenzen in Optimierungsproblemen)
            \item CSP's mit weichen Constraints nennt man \fatsf{Constrained Optimization Problems}
        \end{itemize}
    \end{itemize}
\end{enumerate}
\subsubsection*{How do we solve CSPs?}
Es gibt zwei Hauptmethoden:
\begin{enumerate}
    \item \fatsf{Suche (Search-based CSP Solving)}
    \begin{itemize}
        \item Iterative \fatsf{Zuweisung von Werten} an Variablen
        \item \fatsf{Überprüfung} der Constraints nach jeder Zuweisung
        \item Falls eine \fatsf{Constraint verletzt wird}, erfolgt \fatsf{Backtracking}
        \item Ziel: Eine \fatsf{vollständige konsistente Zuweisung} für alle Variablen
    \end{itemize}
    \item \fatsf{Constraint Propagation}
    \begin{itemize}
        \item Behalten für jede Variable $X_{i}$ eine \fatsf{Menge von möglichen Werten $D_{i}$} bei
        \item Versuchen, \fatsf{unmögliche Werte zu eliminieren}, um die Menge $D_{i}$ zu reduzieren
        \item Beispiel: \fatsf{Arc Consistency (AC-3)} reduziert die Domäne durch Überprüfung binärer Constraints
    \end{itemize}
\end{enumerate}
\subsubsection*{Search-Based CSP Solving}
CSPs können als \fatsf{Suchprobleme} formuliert werden:
\begin{itemize}
    \item \fatsf{Knoten:} Teilweise zugewiesene Variablen
    \item \fatsf{Nachbarn:} Neue Zuweisungen für eine weitere Variable
    \item \fatsf{Startknoten:} Leere Zuweisung
    \item \fatsf{Zielknoten:} Alle Variablen sind zugewiesen und \fatsf{alle Constraints sind erfüllt}
\end{itemize}
\fatsf{Wie bauen wir den Suchraum auf?}
\begin{enumerate}
    \item \fatsf{Start:} Leere Zuweisung
    \item \fatsf{Nachfolger:} Jede Möglichkeit, eine Variable zuzuweisen
    \item \fatsf{Schrittweise Zuweisung weiterer Variablen}
\end{enumerate}
\subsubsection*{Naive Search \& Komplexität}
Die naive Lösung ist eine \fatsf{vollständige Tiefensuche} durch den Suchraum\\
\fatsf{Annahmen:}
\begin{itemize}
    \item n Variablen mit v möglichen Werten pro Variable
    \item Alle Lösungen befinden sich auf \fatsf{Tiefe n} im Suchbaum
\end{itemize}
\fatsf{Berechnung der Komplexität:}
\begin{itemize}
    \item Auf \fatsf{Ebene 1:} n $\times$ v mögliche Zuweisungen
    \item Auf \fatsf{Ebene 2:} $(n - 1) \times v$ mögliche Zuweisungen
    \item Allgemein: \fatsf{Branching Factor auf Tiefe l = (n - l + 1) $\times$ v}
    \item \fatsf{Gesamtanzahl der Blätter im Suchraum:} $n! \times v^{n}$
\end{itemize}
\fatsf{Problem: Exponentielle Komplexität} $\rightarrow$ Sehr ineffizient für große CSPs!
\subsubsection*{Effizientere Lösungen: Commutative Variable Assignments}
\fatsf{Optimierung:} Reihenfolge der Zuweisung ist irrelevant!\\
Beispiel:
\begin{itemize}
    \item Ob wir zuerst WA = rot, NT = grün oder umgekehrt zuweisen, ist egal
    \item \fatsf{Lösung:} Zuweisungen können \fatsf{ohne Reihenfolge-Überprüfung} durchgeführt werden
    \item Dadurch \fatsf{Reduktion der Komplexität} von $n! \times v^{n}$ auf $v^{n}$
\end{itemize}
\fatsf{Effizeinzsteigerung durch Heuristiken:}
\begin{itemize}
    \item \fatsf{Minimum Remaining Values (MRV):} Wähle die Variable mit den \fatsf{wenigsten verbleibenden Optionen} zuerst
    \item \fatsf{Least Constraining Value:} Wähle den Wert, der die \fatsf{wenigsten zukünftigen Einschränkungen verursacht}
\end{itemize} 
\begin{figure}[h]
    \centering
    \includegraphics[width=1\linewidth]{CSP_Images/Commutative Variable Assignments.jpeg}
\end{figure}
\clearpage
\subsubsection*{Backtracking Search in CSPs}
\begin{figure}[h]
    \centering
    \includegraphics[width=1\linewidth]{CSP_Images/Backtracking Search.jpeg}
\end{figure}
Backtracking ist eine \fatsf{tiefensuchbasierte (DFS) Methode}, die schrittweise eine \fatsf{einzelne Variable pro Ebene zuweist}\\
\fatsf{Eigenschaften von Backtracking:}
\begin{itemize}
    \item Grundlegender , uninformativer Suchalgorithmus für CSPs
    \item Fügt eine Variable nach der anderen hinzu, solange keine Konflikte auftreten
    \item Lösung gefunden: Falls eine legale vollständige Zuweisung existiert
    \item Effizienter als brutale Kraft (Naive Suche), aber trotzdem exponentielle Worst-Case-Komplexität
    \item Kann das n-Queens Problem bis n = 25 lösen
    \item Verbesserungen durch Heuristiken möglich
\end{itemize}
\fatsf{Problem von Backtracking:}
\begin{itemize}
    \item Wenn eine inkonsistente Belegung erkannt wird, muss der Algorithmus zur letzten gültigen Zuweisung zurückkehren (Backtracken)
    \item Im Wort-Case exponentielle Laufzeit O($v^{n}$)
    \item Bei großen CSPs zu ineffizient, wenn keine Optimierungstechniken eingesetzt werden
\end{itemize}
\subsubsection*{Constraint Propagation (Beispiel: Sudoku)}
Ein \fatsf{leistungsfähigeres Verfahren} zur Reduzierung der Anzahl möglicher Werte vor der eigentlichen Suche\\
\fatsf{Sudoku als CSP:}
\begin{itemize}
    \item 81 Variablen (Felder des Gitters)
    \item \fatsf{Domäne:} Werte von \fatsf{1 bis 9}
    \item \fatsf{Constraints:}
    \begin{itemize}
        \item Vorab definierte Zahlen sind fest
        \item Jede \fatsf{Zeile, Spalte und 3x3-Box} muss unterschiedlichen Zahlen enthalten (27 Constraints)
    \end{itemize}
\end{itemize}
\fatsf{Idee der Constraint Propagation:}
\begin{itemize}
    \item \fatsf{Schrittweises Eliminieren unmöglicher Werte} für jede Variable
    \item \fatsf{Ziel:} Domäne $D_{i}$ für jede Variable $X_{i}$ zu verkleinern
    \item \fatsf{Beispiel:} Falls ein Feld nur noch eine Möglichkeit hat, muss es diesen Wert annehmen
    \item \fatsf{Effekt:} Reduziert Suchraum drastisch - viele Sudoku-Rätsel lassen sich \fatsf{ohne Backtracking lösen!}
\end{itemize}
\fatsf{Status:} Automatisierte Constraint-Solver lösen selbst schwerste Sudokus extrem schnell!
\subsubsection*{Optimierung von Backtracking: General-Purpose Methods}
\fatsf{Fragen zur Optimierung:}
\begin{enumerate}
    \item Welche Variable sollte als nächstes zugewiesen werden?
    \item In welcher Reihenfolge sollten die Werte getestet werden?
    \item Können wir offensichtliche Sackgassen frühzeitig erkennen?
    \item Können wir von Problemstrukturen profitieren?
\end{enumerate}
\fatsf{Antwort: Intelligente Heuristiken!}
\subsubsection*{Heuristics for CSPs}
\begin{enumerate}
    \item \fatsf{Minimum Remaining Values (MRV) Heuristic}
    \begin{itemize}
        \item \fatsf{Wähle die Variable mit den wenigsten verbleibenden konsistenten Werten}
        \item Grund: Diese Variable wird am schwierigsten zuzuweisen sein - also priorisieren!
        \item \fatsf{Effekt:} Reduziert frühzeitig die Anzahl der Möglichkeiten und verhindert tiefe Fehlersuche
    \end{itemize}
    \fatsf{Bespiel:} Wenn ein Sudoku-Feld nur noch \fatsf{eine mögliche Zahl} hat, setzen wir es direkt!
    \item \fatsf{Degree Heuristic}
    \begin{itemize}
        \item \fatsf{Wähle die Variable mit den meisten Constraints auf verbleibenden Variablen}
        \item Grund: Diese Variable schränkt andere am stärksten ein - je früher wir sie setzen, desto besser
        \item \fatsf{Effekt:} Verringert die Unsicherheiten für verbleibende Variablen
    \end{itemize}
    \fatsf{Beispiel:} Beim \fatsf{Vier-Farben-Problem} wählen wir zuerst das Gebiet mit den meisten Nachbarn, da es die anderen stakt beeinflusst
    \begin{figure}[h]
        \centering
        
\includegraphics[width=1\linewidth]{CSP_Images/Degree Heuristic.jpeg}
    \end{figure}
    \item \fatsf{Least Constraining Value (LCV) Heuristic}
    \begin{itemize}
        \item \fatsf{Für eine Variable: Wähle den Wert, der die wenigsten Werte für andere Variablen blockiert}
        \item Grund: Dadurch bleibt mehr Flexibilität für spätere Zuweisungen
        \item \fatsf{Effekt:} Verhindert Sackgassen und unnötige Backtracking-Schritte
    \end{itemize}
    \fatsf{Beispiel:} WEnn wir eine Farbe für ein Gebiet in einer Landkarte wählen, nehmen wir \fatsf{diejenige, die möglichst viele Optionen für die Nachbargebiete übrig lässt}
    \begin{figure}[h]
        \centering
        \includegraphics[width=1\linewidth]{CSP_Images/LCV.jpeg}
    \end{figure}
\end{enumerate}
\fatsf{Kombinierte Anwendung der Heuristiken}\\
\fatsf{MRV $\rightarrow$ Degree $\rightarrow$ LCV} ergibt eine extrem effiziente Suchstrategie!
\begin{enumerate}
    \item \fatsf{MRV:} Wähle schwierigste Variable zuerst
    \item \fatsf{Degree:} Falls mehrere Optionen, wähle die mit den meisten Constraints
    \item \fatsf{LCV:} Falls mehrere Werte möglich, wähle den flexibelsten
\end{enumerate}
\subsection{Constraint Propagation}
\fatsf{Idee:}
\begin{itemize}
    \item Reduziert frühzeitig die Anzahl möglicher Werte für Variablen
    \item Erlaubt es, \fatsf{Suche effizienter zu gestalten} (z.B. bei Sudoku)
    \item Ziel: \fatsf{Wiedersprüche frühzeitig erkennen} und unnötige Berechnungen vermeiden
\end{itemize}
\subsubsection*{Node Consistency (1-Consistency)}
Ein \fatsf{Knoten (Variable) ist konsistent}, wenn alle seine möglichen Werte mit \fatsf{unären Constraints} übereinstimmen\\
Beispiel: \fatsf{Sudoku}
\begin{itemize}
    \item Varab ausgefüllte Zahlen sind automatisch \fatsf{node-consistent}, weil sie auf einen einzigen Wert beschränkt sind
\end{itemize}
\subsubsection*{Arc Consistency (2-Consistency)}
Ein \fatsf{Bogen (Arc) ist konsistent}, wenn für \fatsf{jeden Wert $X_{i} in der Domäne$} eine gültige Wahl für $X_{j}$ existiert, sodass die \fatsf{Constraint-Bedingung zwischen den beiden Variablen erfüllt ist}
\begin{itemize}
    \item Prüft \fatsf{binäre Constraints} (zwei Variablen)
    \item Falls eine Variable \fatsf{keine möglichen Werte mehr hat}, gibt es \fatsf{keine Lösung}
    \item In der Praxis häufig genutzt
\end{itemize}
\fatsf{Besipiel:} In \fatsf{Vier-Farben-Problem} bedeutet Arc Consistency, dass für jede Farbe einer Region mindestens eine zulässige Nachbarfarbe existiert
\subsubsection*{Forward Checking}
\fatsf{Idee:}\\
Nach jeder Variablenzuweisung werden die \fatsf{Domänen der noch nicht zugewiesenen Variablen reduziert}
\begin{itemize}
    \item Falls eine Variable \fatsf{keine möglichen Werte mehr hat $\rightarrow$ Sofortiger Backtrack!}
    \item Besonders hilfreich, um \fatsf{Fehlschläge frühzeitig zu erkennen}
\end{itemize}
\fatsf{Beispiel:} Im \fatsf{Sudoku:} Falls wir eine 3 in eine Zeile setzen, entfernen wir 3 aus allen möglichen Werte der zugehörigen Zeile, Spalte und Box
\begin{figure}[h]
    \centering
    \includegraphics[width=1\linewidth]{CSP_Images/Forward Checking.jpeg}
\end{figure}
\subsubsection*{AC-3 Algorithmus}
Ein effizienter Algorithmus zur \fatsf{Herstellung von Arc Consistency}
\fatsf{Prinzip:}
\begin{itemize}
    \item Überprüft systematisch, ob für jede Variable mindestens ein passender Wert für ihre Nachbarn existiert
    \item Falls eine Änderung vorgenommen wird, werden betroffene Bögen erneut überprüft
\end{itemize}
\fatsf{Laufzeit:} O($ed^{2}$) (e = Anzahl der Bögen, d = maximale Domänengröße)\\
\fatsf{Beispiel:} Im \fatsf{Sudoku} reduziert AC-3 die Anzahl möglicher Werte drastisch und kann viele Rätsel \fatsf{ohne Backtracking lösen!}
\subsubsection*{Path Consistency (3-Consistency)}
\fatsf{Idee:}
\begin{itemize}
    \item Überprüft, ob \fatsf{zwei Variablen konsistent mit einer dritten Variable $X_{m}$ sind}
    \item Erlaubt feinere Reduktion des Suchraums als Arc Consistency
\end{itemize}
\fatsf{Formale Definition:}\\
Ein Paar ($X_{i}, X_{j}$) ist path-consistent mit $X_{m}$, wenn für jede gültige Zuweisung von $X_{i}$ und $X_{j}$ eine gültige Wahl für $X_{m}$ existiert\\
\fatsf{Beispiel:} Im \fatsf{n-Queens Problem} wird überprüft, ob zwei Damen nicht nur einander, sondern auch weitere Damen in einer dritten Dimension (Diagonale) beeinflussen
\subsubsection*{k-Consistency}
Generalisierung von Konsistenz auf \fatsf{mehr als zwei Variablen}
\begin{itemize}
    \item \fatsf{1-Consistency:} Node Consistency
    \item \fatsf{2-Consistency:} Arc Consistency
    \item \fatsf{3-Consistency:} Path Consistency
    \item \fatsf{k-Consistency:} Ein Set von k Variablen muss konsistent sein
\end{itemize}
\fatsf{Problem:}
\begin{itemize}
    \item Überprüfung von \fatsf{k-Consistency ist exponentiell in k}
    \item Daher wird \fatsf{Arc Consistency (2-Consistency) am häufigsten verwendet}
\end{itemize}
\subsubsection*{Constraint Propagation \& Backtracking Search}
\fatsf{Idee:} Kombination von Constraint Propagation mit Backtracking, um den Suchraum effizeint zu reduzieren\\
\fatsf{Prinzip:}
\begin{enumerate}
    \item \fatsf{Nach jeder Variablenzuweisung} wird ein Constraint Propagation Algorithmus angewendet
    \item Dadurch \fatsf{reduziert sich die Anzahl möglicher Werte für andere Variablen}
    \item Dies führt zu \fatsf{weniger Backtracking} und verbessert die Effizienz
\end{enumerate}
\subsubsection*{Mögliche Algorithmen zur Constraint Propagation}
\begin{itemize}
    \item \fatsf{Forward Checking}
    \begin{itemize}
        \item Entfernt direkt nach jeder Zuweisung unzulässige Werte aus den verbleibenden Variablen
        \item Falls eine Variable \fatsf{keine Werte mehr hat $\rightarrow$ Backtrack}
    \end{itemize}
    \fatsf{Beispiel:} In \fatsf{Sudoku} werden nach jeder Zahleneitragung die Domänen der betroffenen Felder reduziert
    \item \fatsf{AC-3 Algorithmus}
    \begin{itemize}
        \item Stellt \fatsf{Arc Consistency} für alle Bögen her
        \item Falls eine Variable \fatsf{keine Werte mehr hat $\rightarrow$ Backtrack}
    \end{itemize}
    \fatsf{Beispiel:} Kann oft CSPs \fatsf{ohne Backtracking vollständig lösen!}
\end{itemize}
\subsubsection*{Local Search für CSPs}
\fatsf{Modifikationen für CSP:}
\begin{itemize}
    \item Arbeitet mit \fatsf{kompletten Zuständen} (alle Variablen sind belegt)
    \item Erlaubt \fatsf{Zustände mit verletzten Constraints}
    \item \fatsf{Operatoren:} Weisen Variablen \fatsf{neue Werte zu}, um Konflikte zu minimieren
\end{itemize}
\fatsf{Min-Conflicts Heuristic}
\begin{enumerate}
    \item Wähle eine zufällige Variable mit \fatsf{verletzt} Constraints
    \item Wechsle ihren Wert auf den, der \fatsf{die wenigsten Konflikte verursacht}
    \item Dies entspricht \fatsf{Hill Climbing mit h(n) = Anzahl der verletzten Constraints}
\end{enumerate}
\fatsf{Beispiel:} Sehr effizient für \fatsf{n-Queens Problem} (löst 1 Mio. Queens in Sekunden!)
\subsubsection*{Problem-Dekomposition}
\fatsf{Idee:}
\begin{itemize}
    \item \fatsf{Große CSPs in kleinere Teilprobleme zerlegen} $\rightarrow$ drastische Reduzierung der Komplexität!
\end{itemize}
\begin{itemize}
    \item \fatsf{Standard-Komplexität:} $O(d^{n})$ für n Variablen mit d möglichen Werten
    \item \fatsf{Nach Dekomposition:} Falls Teilprobleme je c Variablen haben, reduziert sich die Komplexität auf $O(n/c \times d^{c})$
\end{itemize}
\fatsf{Beispiel:} Für n = 80, d = 2, c = 20 ergibt sich eine massive \fatsf{Verkleinerung des Suchraums!} $\rightarrow$ $2^{80} \neq 4 \times 2^{20}$
\subsubsection*{Tree-structured CSPs}
\fatsf{Theorem:}
\begin{itemize}
    \item Jedes CSP mit Baumstruktur kann in linearen Zeit gelöst werden!
    \item Komplexität: O($n \times d^{2}$)
\end{itemize}
\fatsf{Algorithmus:}
\begin{enumerate}
    \item Wähle eine Variable als \fatsf{Wurzel}
    \item Ordne alle Variablen so, dass \fatsf{Eltern vor ihren Kinderns kommen}
    \item \fatsf{Von den Blättern bis zur Wurzel}
    \begin{itemize}
        \item Mache jeden Bogen ($X_{i}, X_{j}$) \fatsf{arc-consistent}
    \end{itemize}
    \item \fatsf{Von der Wurzel abwärts:}
    \begin{itemize}
        \item Weise jeder Variablen eine \fatsf{gültige Belegung} zu
    \end{itemize}
\end{enumerate}
\fatsf{Beispiel:} Falls der CSP-Graph ein \fatsf{Baum} ist, kann er \fatsf{ohne Backtracking} gelöst werden!
\begin{figure}[h]
    \centering
    \includegraphics[width=1\linewidth]{CSP_Images/Tree-structured CSPs.jpeg}
\end{figure}
\subsubsection*{Nearly Tree-structured CSPs}
\fatsf{Problem:} Baumstruktur ist selten in realten Problemen (oft gibt es Zyklen)\\
\fatsf{Zwei Techniken zur Reduktion auf Bäume:}
\begin{enumerate}
    \item \fatsf{Cutset Conditioning}
    \begin{itemize}
        \item Entferne eine \fatsf{kleine Anzahl von Variablen}, sodass der Graph \fatsf{baumförmig} wird
        \item Löse dann das Problem mit \fatsf{Baum-Algorithmus}
    \end{itemize}
    \item \fatsf{Knoten-Zusammenfassung}
    \begin{itemize}
        \item \fatsf{Gruppiere Variablen} zu \fatsf{Teilproblemen mit Baumstrutkur}
        \item Löse die kleineren Teilproblem \fatsf{unabhängig}
    \end{itemize}
\end{enumerate}
\fatsf{Beispiel für Cutset Conditioning:} Entferne \fatsf{einige Variablen}, sodass der Rest \fatsf{baumartig} wird $\rightarrow$ dann effizient lösbar!
\begin{figure}[h]
    \centering
    \includegraphics[width=1\linewidth]{CSP_Images/Cutset Conditioning.jpeg}
\end{figure}